\section{Finite randomness, counting, circuits, and lower bounds}
\label{sec:random-lower}

The central objects in this section are finite tables and finite computation
traces. A probability is a rational abbreviation for a ratio of two natural
numbers; every inequality below has an explicitly stated integer form.
The standard complexity results provide the targets
\cite{adleman1978,lautemann1983,shannon1949,katz2011}.
The proofs here reconstruct them with bounded searches, positive correctness
statements, and numerical induction invariants. In particular, the assertion
that a finite bad set is smaller than its ambient set produces its least
missing element. It is not an appeal to an unbounded choice principle.

\subsection{The counting primitives and the randomized presentations}

For a coded Boolean table or a fixed primitive-recursive Boolean predicate
$E(i,\vec a)$, put
\[
 C_E(N,\vec a)=\sum_{i<N}[E(i,\vec a)=1].
\]
The brackets have values zero and one. The defining recursion is
$C_E(0)=0$ and $C_E(N+1)=C_E(N)+[E(N)=1]$.
Sets in this section are abbreviations for such predicates on bounded
intervals; their sizes mean these sums. Tuples are encoded in a fixed radix,
so a tuple of $k$ elements below $Q$ is a number below $Q^k$.
The intended language is always specified by a Boolean function $L(x)$,
or by a fixed bounded definition already compiled to such a function.

\begin{lemma}[R01: finite counting calculus]\label{thm:R01}
The following statements have primitive-recursive witnesses and proofs by
quantifier-free induction in the bounded compilation described earlier.
\begin{enumerate}
\item $0\leq C_E(N)\leq N$, and $C_E(N)>0$ if and only if there is an
      $i<N$ with $E(i)=1$. If the count is positive, the least such $i$ is
      obtained by a bounded scan.
\item If $E(i)\leq F(i)$ for all $i<N$, then $C_E(N)\leq C_F(N)$.
      For a finite family $(E_j)_{j<k}$,
      $|\bigcup_{j<k}E_j|\leq\sum_{j<k}|E_j|$.
\item $|E\times F|=|E||F|$. More generally, a sum of products over a
      finite Cartesian product factors into the product of the sums.
\item If a permutation of $[0,N)$ carries $E$ onto $F$, their counts agree.
      If an explicitly described involution without fixed points exchanges
      $E$ and its complement, then $2|E|=N$.
\item If $|E|<N$, the least $i<N$ outside $E$ exists and is found by bounded
      search. The same assertion applies to a union whose displayed upper
      bound is strictly below $N$.
\end{enumerate}
\end{lemma}
\begin{proof}
For the first two assertions, add the next zero-or-one term to both sides
of the maintained inequality. For the union bound use the pointwise
inequality $[\bigvee_{j<k}E_j(i)]\leq\sum_{j<k}[E_j(i)]$, then sum in $i$.
For the product assertion, split the radix enumeration into consecutive
blocks of the second radix and apply distributivity; induction in the
number of blocks proves the equality. A permutation rearranges a finite
sum, as is proved by tracking the counts of its already visited indices.
An exchanging involution is a bijection between the two parts, so their
counts add to $N$ and are equal. Finally the scan for a first missing
element has invariant: either an already recorded index is missing, or
every earlier index belongs to $E$. If the scan ended with no missing
index, its count would be $N$, contradicting the strict bound. All these
invariants contain only bounded quantifiers over explicitly coded tables.
Their Boolean compilations, rather than induction on arbitrary formulas,
are the object-theory induction instances.
\end{proof}

A randomized polynomial-time presentation consists of a fixed deterministic
machine $M$, an explicit polynomial clock, and a polynomial $r(n)$ bounding
the number of independent random bits. Unused bits are ignored, so each
input of length $n$ has exactly
\[
 Q(n)=2^{r(n)}
\]
equally weighted seeds. Its clocked output is denoted $M(x,u)\in\{0,1\}$,
for $u<Q(n)$. Define
\[
 a_M(x)=\sum_{u<Q(n)}M(x,u),\qquad
 e_{M,L}(x)=\sum_{u<Q(n)}[M(x,u)\ne L(x)].
\]
All bounds are on binary input length; enumeration of the $Q(n)$ seeds is
used in the proof theory, not asserted to be polynomial time.

An RP presentation satisfies
\[
 L(x)=0\Longrightarrow a_M(x)=0,\qquad
 L(x)=1\Longrightarrow 2a_M(x)\geq Q(n).
\]
A coRP presentation has $a_M(x)=Q(n)$ on yes inputs and
$2a_M(x)\leq Q(n)$ on no inputs. A BPP presentation satisfies
$3e_{M,L}(x)\leq Q(n)$ for every $x$. A PP presentation satisfies the
positive equivalence
\[
 L(x)=1\quad\Longleftrightarrow\quad 2a_M(x)>Q(n).
\]
The left side is a decidable Boolean equality. These conventions prevent
an undecidable proposition from being inserted as an unspecified
``intended language'' behind two one-sided implications.

\begin{theorem}[R02: one-sided randomness and witnesses]\label{thm:R02}
$\mathrm{P}\subseteq\mathrm{RP}\subseteq\mathrm{NP}$ and
$\mathrm{P}\subseteq\mathrm{coRP}\subseteq\mathrm{coNP}$.
An RP computation repeated independently $k\geq1$ times, accepting if any
trial accepts, still has no false positives and has at most $Q^k/2^k$
false-negative seeds. The analogous AND construction amplifies coRP.
In particular RP and coRP are contained in BPP.
\end{theorem}
\begin{proof}
A deterministic decider may ignore its random seed. On an RP yes input,
$2a_M\geq Q>0$ gives $a_M>0$; Lemma~\ref{thm:R01} supplies an accepting
seed, of polynomial length, as an NP witness. Conversely an accepting
seed is impossible when $L(x)=0$. Complementing both the output and the
language gives the coRP assertion.

For amplification let $b=Q-a_M$ on a yes input, so $2b\leq Q$. Exactly
$b^k$ seed tuples make all $k$ trials reject, by the product identity.
Thus $2^k b^k\leq Q^k$. On a no input each trial rejects, so the OR still
rejects on every tuple. Two trials give error at most $1/4$, which is at
most $1/3$. AND amplification follows by exchanging zero and one.
The machine uses $kr(n)$ random bits and $k$ times the original clock,
plus polynomial overhead for counters and the accumulated Boolean result.
\end{proof}

\begin{theorem}[R03: integer BPP amplification]\label{thm:R03}
For every $d\geq1$, taking the majority of $12d$ independent BPP trials
(breaking ties in a fixed way) gives
\[
 2^d e_{\mathrm{amp},L}(x)<Q(n)^{12d}.
\]
The number of random bits is $12d\,r(n)$; if $d$ is polynomial in $n$,
the amplified machine is polynomial time.
\end{theorem}
\begin{proof}
Fix $x$, let $b=e_{M,L}(x)$, and put $k=2t$. For a seed tuple $\vec u$,
let $j(\vec u)$ be its number of erroneous trials. Factorization of finite
sums gives the exact identity
\[
 \sum_{\vec u\in[0,Q)^k}2^{j(\vec u)}=(Q+b)^k.
\]
Indeed each coordinate contributes weight two on its $b$ erroneous seeds
and weight one on its $Q-b$ correct seeds, whose sum is $Q+b$.
An erroneous majority has $j(\vec u)\geq t$, regardless of the tie rule.
If $E$ is its error count, then
\[
 2^tE\leq(Q+b)^{2t},\qquad
 3^{2t}2^tE\leq(4Q)^{2t},
\]
because $3b\leq Q$. Cancelling the positive factor $2^t$ yields
$9^tE\leq8^tQ^{2t}$. Now choose $t=6d$. The explicit integer check
$2\cdot8^6=524288<531441=9^6$, followed by induction on $d$, gives
$2^d8^{6d}<9^{6d}$. Combining the inequalities proves the result.
No logarithm, real exponential, Chernoff bound, or integration is used.
\end{proof}

\begin{theorem}[R04: elementary closure of randomized classes]\label{thm:R04}
BPP is closed under complement, union, intersection, and polynomial-time
many-one preimages. RP and coRP are each closed under union, intersection,
and polynomial-time many-one preimages.
\end{theorem}
\begin{proof}
For complement, replace $M$ by $1-M$: its error indicator relative to
$1-L$ is identical. For BPP union and intersection, first amplify both
machines with $d=3$. Each has error less than $1/8$. Run them with separate
seeds and take respectively OR or AND. A wrong combined answer requires
at least one wrong component answer, hence the union bound gives total
error less than $1/4<1/3$ after clearing the seed denominators.

For RP union, run both original machines and take OR. There are no false
positives, and on a yes input at least one component succeeds on at least
half its seeds. For RP intersection first amplify both to failure at most
$1/4$, and take AND. On a yes input, at most half the seed pairs fail by
the union bound; on a no input one component always rejects. Complement
these constructions to obtain the two coRP closures. Finally compute a
reduction $f(x)$ deterministically and run the chosen randomized machine
on $f(x)$. Polynomial output length composes with its clock and random-bit
bound. Its count inequality is exactly the old inequality for $f(x)$.
\end{proof}

\subsection{Counting functions, majority computation, and advice}

A $\#\mathrm P$ presentation is a polynomial-time Boolean verifier
$V(x,w)$ and a polynomial witness length $p(n)$, with
\[
 f(x)=\sum_{w<2^{p(n)}}V(x,w).
\]
This is an actual primitive-recursive natural-valued function. Witnesses
are bit strings of an exact declared length. A convention that merely
allows arbitrary trailing padding would change counts and is therefore
not silently interchangeable with this one.

\begin{theorem}[R05: algebra of witness counts]\label{thm:R05}
$\#\mathrm P$ is closed under addition, multiplication, and polynomially
bounded sums of uniformly polynomial-time verifiers. Its positive supports
are precisely NP. Every nonnegative polynomial-time computable function
of polynomial output length is a $\#\mathrm P$ function. A fixed
presentation $f(x)=\#\{w<2^{p(n)}:V(x,w)=1\}$ also gives a
$\#\mathrm P$ presentation of $2^{p(n)}-f(x)$.
\end{theorem}
\begin{proof}
For addition, pad both witness lengths to their maximum by requiring the
new padding bits to be zero. Prepend a tag bit selecting the first or
second verifier. There is exactly one padded witness for each original
witness, so the tag partitions the accepted witnesses into two disjoint
sets of sizes $f(x)$ and $g(x)$. For multiplication concatenate one witness
for each verifier and accept exactly when both verifiers accept. The
product identity in Lemma~\ref{thm:R01} gives count $f(x)g(x)$.

For a sum over $i<k(n)$, use an index of length
$\lceil\log_2(k(n)+1)\rceil$ and a common polynomial witness length;
reject invalid indices and require unused witness padding to be zero.
The accepted pairs partition by $i$. A positive witness count is
equivalent, with a bounded least-witness map, to existence of an accepted
witness, which is exactly the NP presentation. If $f(x)<2^{p(n)}$ is
computed in polynomial time, accept $w$ precisely when $w<f(x)$; its
count is $f(x)$. Finally complement the Boolean verifier to obtain the
count $2^{p(n)}-f(x)$. This last result is relative to a specified witness
cube, and does not assert closure under arbitrary subtraction.
\end{proof}

\begin{theorem}[R06: PP containments and complement]\label{thm:R06}
$\mathrm{NP}\cup\mathrm{coNP}\cup\mathrm{BPP}\subseteq\mathrm{PP}
\subseteq\mathrm{PSPACE}$, and PP is closed under complement. Every
$\#\mathrm P$ function is computable using polynomial work space and
exponential time.
\end{theorem}
\begin{proof}
For an NP verifier with $Q=2^{p(n)}$ witnesses and acceptance count $a$,
use a tag bit and $p(n)$ other bits. On tag zero accept; on tag one run
the verifier. Of the $2Q$ seeds, exactly $Q+a$ accept, which is more than
half precisely when $a>0$. Thus the constructed PP language is exactly
the NP support, with both implications witnessed by finite arithmetic.

For complement closure, first pad the random length of a PP machine so
it is at least one. Write $Q=2H$. Use $2Q$ seeds $(b,u)$ with $u<Q$.
For $b=0$ accept if the original machine rejects. For $b=1$ accept if
$u<H+1$. The new acceptance count is
\[
 Q-a+H+1=3H+1-a.
\]
It is strictly greater than half of $2Q$, namely $Q=2H$, precisely when
$a\leq H$. This is the complement of $a>H$. The construction fits the
seed space since $H\geq1$ and $H+1\leq Q$. It also gives
$\mathrm{coNP}\subseteq\mathrm{PP}$ by the preceding NP construction.
A BPP yes input has $3a\geq2Q$, whence $2a>Q$; a no input has
$3a\leq Q$, whence $2a\leq Q$. Its original machine therefore already
satisfies the PP specification.

To compute a count, enumerate all $p(n)$-bit witnesses, retaining the
current witness, a $p(n)+1$-bit accumulator, and the work tape of the
polynomial-time verifier. Reuse the work tape between witnesses. The
workspace is polynomial and the time is $2^{p(n)}$ times a polynomial.
Comparing the resulting count to the half-cube threshold proves the PP
containment. The invariant after witness $w$ is that the accumulator
equals the sum on earlier witnesses; its Boolean form is bounded.
\end{proof}

\begin{theorem}[R07: BPP has polynomial advice and circuits]\label{thm:R07}
From a BPP presentation one can construct, for every length $n$, a
polynomial-length advice string $a_n$ that makes a fixed polynomial-time
machine correct on all $n$-bit inputs. The function $n\mapsto a_n$ can be
chosen primitive recursively. Consequently
$\mathrm{BPP}\subseteq\mathrm P/\mathrm{poly}$ in the finite circuit
and advice readings.
\end{theorem}
\begin{proof}
Amplify with $d=n+2$. Let $m(n)=12(n+2)r(n)$ and let $Q'=2^{m(n)}$ be
the amplified seed count. For each $n$-bit input $x$, fewer than
$Q'/2^{n+2}$ seeds give a wrong answer. The union over the $2^n$ inputs
has fewer than $Q'/4$ members, after multiplying through by $2^{n+2}$.
It is therefore not the whole seed cube. Choose the least seed outside
this union and call it $a_n$. The search predicate is a conjunction over
$2^n$ inputs of finite machine computations and comparisons with $L(x)$.
Thus both the search bound $Q'$ and its Boolean predicate are primitive
recursive. No countable choice of unrelated advice strings is necessary.

The deterministic advice machine runs the amplified machine with seed
$a_n$. Its clock is polynomial and its output equals $L(x)$ for every
input of length $n$, by the definition of the selected seed. Unrolling
this fixed clocked computation into a Boolean circuit, with the advice
bits hardwired, produces a polynomial-size circuit. Conversely a circuit
description of polynomial length can be used as advice by a universal
circuit evaluator; the gates are evaluated in topological order in
polynomial time. This explains the equivalence of the two formulations
used here. It does not say that the advice can be found in polynomial
time. Nor does it say that arbitrary languages in $\mathrm P/\mathrm{poly}$
are recursive: the primitive-recursive advice selector is available here
because it is constructed from the given BPP presentation.
\end{proof}

\subsection{A fully finite proof that BPP lies at the second hierarchy level}

\begin{theorem}[R08: the translate-covering construction]\label{thm:R08}
Every BPP presentation has both a $\Sigma^\mathrm P_2$ presentation and a
$\Pi^\mathrm P_2$ presentation, with explicit polynomial bounds and
primitive-recursive witness searches. Hence
\[
 \mathrm{BPP}\subseteq\Sigma^\mathrm P_2\cap\Pi^\mathrm P_2.
\]
\end{theorem}
\begin{proof}
Let $r=r(n)$ be the original random-bit bound and set
\[
 d=16(r+1)^2,\qquad m=12dr,\qquad Q=2^m,\qquad k=m+1.
\]
Amplify as in Theorem~\ref{thm:R03}; call the resulting predicate $A(x,u)$.
On a yes input its rejecting set has size $b$ with $2^d b<Q$; on a no
input its accepting set has size $a$ with $2^d a<Q$.
We require $k<2^d$. To verify it without an asymptotic exception, put
$v=r+1\geq1$. Then $k\leq193v^3<2^{8+3v}\leq2^{16v^2}$, using
$v\leq2^v$ and $8+3v\leq16v^2$. All constants and inequalities are
integer ones.

Bitwise exclusive-or on $m$-bit words is denoted $\oplus$; translation
by a fixed word is a permutation. Consider
\[
 \exists s_0,\ldots,s_{k-1}<Q\quad
 \forall u<Q\quad \bigvee_{i<k} A(x,u\oplus s_i)=1.                 \tag{R8a}
\]
Its witness block has length $km$, its universal block has length $m$,
and its matrix is decidable by $k$ polynomial-time runs. It is therefore
a $\Sigma^\mathrm P_2$ formula.

On a yes input, fix $u$. Exactly $b^k$ of the $Q^k$ shift tuples have all
$u\oplus s_i$ in the rejecting set: each coordinate is a permutation of
the seed cube. At most $Qb^k$ shift tuples leave some $u$ uncovered, by
the union bound. Since $2b<Q$ and $k=m+1$,
\[
 Qb^k < Q\,Q^k/2^k = Q^k/2 < Q^k.
\]
Thus a least covering tuple exists. This proves (R8a) constructively.
On a no input, any $k$ translates of the accepting set have union size
at most $ka<Q$, since $k<2^d$ and $2^d a<Q$. The least uncovered $u$
refutes that tuple. This supplies the reverse implication with a bounded
counterexample function.

For completeness, the dual positive presentation is
\[
 \forall s_0,\ldots,s_{k-1}<Q\quad
 \exists u<Q\quad \bigwedge_{i<k} A(x,u\oplus s_i)=1.              \tag{R8b}
\]
On a yes input the $k$ translates of the small rejecting set cannot cover
the cube, so their least uncovered point supplies $u$. On a no input the
large rejecting set has a covering tuple by the already proved argument,
and that tuple refutes (R8b). This proves the $\Pi^\mathrm P_2$ reading
directly, without using unrestricted classical contraposition. The special
case $m=0$ is harmless: $Q=k=1$, and the strict error bounds force zero
errors. Every search in the argument is bounded by $Q$ or $Q^k$.
\end{proof}

\subsection{Zero-error computation through finite truncations}

A bounded zero-error trial returns one of $0,1,?$ in polynomial time.
It has \emph{trace soundness} if every definite answer equals $L(x)$,
and \emph{success density at least one half} if at least half its seeds
produce a definite answer. These are finite Boolean and counting tests.
For a possibly longer-running random machine $Z$, normalize it to inspect
at most one fresh random bit per step. For a prefix of $T$ bits, let
$c_T(x,u)$ be the smaller of its halting time and $T$, defined by a
$T$-step simulation. The finite expected-time certificate with bound $p$
is the arithmetical statement
\[
 \forall x\ \forall T\qquad
 \sum_{u<2^T}c_T(x,u)\leq p(|x|)2^T.                              \tag{R9}
\]
The certificate is accompanied by trace soundness at every truncation.
This is the precise zero-error expected-time reading used here.

\begin{theorem}[R09: zero-error equivalence and restart bounds]\label{thm:R09}
The following presentations are interconvertible by explicit machine
constructions: (i) a bounded zero-error trial of success density at least
one half; (ii) an RP presentation and a coRP presentation of the same
Boolean language; (iii) a trace-sound random machine with a polynomial
finite expected-time certificate (R9). Thus the finite arithmetical
reading of $\mathrm{ZPP}=\mathrm{RP}\cap\mathrm{coRP}$ is provable in
\ipra. Repeating (i) for at most $j$ rounds leaves failure density at most
$2^{-j}$, and its average number of rounds is at most two for every cap.
\end{theorem}
\begin{proof}
From (i), accept only on output one to get RP. Accept on output one or
question mark to get coRP. Trace soundness eliminates the forbidden
one-sided errors, and the definite-output density gives the required
one-half success inequalities.

Conversely let $R$ be an RP test for $L$ and let $S$ be an RP test for
$1-L$, obtained by complementing a coRP test. Run them using disjoint
seeds. Return one if $R$ accepts; otherwise return zero if $S$ accepts;
otherwise return question mark. On a yes input, $S$ never accepts and
$R$ accepts on at least half its seeds. On a no input the roles reverse.
Both a definite wrong answer and a conflict of two valid answers are
therefore excluded by trace soundness. This constructs (i).

Pad one trial to a fixed polynomial duration $p(n)\geq1$, so that
independent blocks of $p(n)$ fresh bits suffice. If $Q=2^{p(n)}$ and
$b\leq Q/2$ is the number of inconclusive blocks, exactly $b^j$ block
tuples fail all $j$ trials. Hence $2^j b^j\leq Q^j$. If $N_j$ is the
number of rounds used by the procedure capped at $j$, finite double
counting gives
\[
 \sum_{\vec u\in[0,Q)^j} N_j(\vec u)
 =\sum_{i=0}^{j-1} b^i Q^{j-i}
 \leq Q^j\sum_{i=0}^{j-1}2^{-i}\leq2Q^j.
\]
The displayed fractions abbreviate the equivalent inequality after
multiplication by $2^{j-1}$ when $j>0$; the case $j=0$ is zero.
The equality counts a round whenever all preceding rounds were
inconclusive. Running successive independent trials until one succeeds
therefore gives (R9) with bound $2p(n)$ at every finite step cap: charge
at most $p(n)$ steps for each entered round, and apply the same finite
tail sum, including the possibly incomplete last block. More explicitly, let
$j=\lceil T/p(n)\rceil$ and extend each $T$-bit prefix to $jp(n)$ bits.
For every extension, $c_T\leq p(n)N_j$. Summing over all extensions and
using their equal multiplicity $2^{jp(n)-T}$ gives
$2^{jp(n)-T}\sum_{u<2^T}c_T(x,u)\leq2p(n)2^{jp(n)}$; cancelling
the positive multiplicity gives (R9). The case $T=0$ is immediate. The trace remains
sound because any eventual answer came from a sound trial. This proves
(i)$\Rightarrow$(iii) without asserting termination on every infinite
random tape.

Finally assume (iii), increase $p$ to at least one if necessary, and cap
the machine at $T=2p(n)$ steps. If $B$ prefixes have not halted by that
cap, each contributes exactly $T$ to the sum in (R9). Thus
$TB\leq p(n)2^T$, and cancellation of $p(n)>0$ yields
$2B\leq2^T$. Return question mark on these prefixes and the observed
answer otherwise. This is (i). The argument is the finite counting form
of Markov's inequality; no measure or assertion of universal random-tape
termination enters it.
\end{proof}


\subsection{Circuit counting with an explicit hard family}

For this subsection, a circuit on $n$ input bits is a finite straight-line
program. It has input wires, two constant wires, and topologically ordered
gates. Each gate applies one of the sixteen binary Boolean functions to two
earlier wires. The output is a specified wire. This finite basis contains
AND, OR, NOT, and copying; using it instead of a smaller complete basis
changes the relevant size bounds only by fixed factors. A truth table is
an integer $h<2^{2^n}$, whose $x$th bit is the value on input $x<2^n$.

\begin{theorem}[R10: quantitative circuit counting]\label{thm:R10}
The number of $n$-input Boolean functions computable with at most $s$ gates
is at most
\[
 B(n,s)=(n+s+2)\bigl(16(n+s+2)^2\bigr)^s.
\]
If $B(n,s)<2^{2^n}$, a primitive-recursive bounded search produces a truth
table not computed by any such circuit.
\end{theorem}
\begin{proof}
Pad a shorter circuit to exactly $s$ gates with unused copy gates, retaining
its output wire. For an upper bound, allow every gate to choose its two
sources from all $n+s+2$ wire positions, even the positions that would be
invalid because they are not earlier. There are at most
$16(n+s+2)^2$ descriptions for each gate and at most $n+s+2$ choices of
output. Every valid padded circuit occurs in this bounded enumeration,
and different descriptions may compute the same function. Consequently
the number of computed truth tables is bounded by $B(n,s)$.

Enumerate a circuit's truth table by evaluating its gates in order on
each of the $2^n$ inputs. Invalid descriptions are marked invalid; they
are not used as extra circuits. Scan the $2^{2^n}$ truth tables and choose
the first one distinct from the tables of all valid descriptions. If no
such table existed, marking the computed tables in an initially zero
bit vector would mark all $2^{2^n}$ positions with at most $B(n,s)$ marks.
The invariant that the number of marked positions is at most the number
of processed descriptions contradicts the strict inequality. The scan,
evaluation, marks, and equality comparisons are all bounded and primitive
recursive. They provide the witness, not merely a double-negated existence.
\end{proof}

\begin{theorem}[R11: a concrete Shannon bound]\label{thm:R11}
For $n\geq4$ let
\[
 s(n)=\left\lfloor\frac{2^n}{16(n+1)}\right\rfloor.
\]
At most $2^{2^{n-1}}$ truth tables on $n$ bits have circuits of size at most
$s(n)$. In particular a least hard table $h_n$ exists, and at least
$2^{2^n}-2^{2^{n-1}}$ tables require more gates.
\end{theorem}
\begin{proof}
Put $s=s(n)$. Since $s\leq2^{n-4}$ and
$n+2+2^{n-4}\leq2^n$ for $n\geq4$, we have $n+s+2\leq2^n$.
The latter elementary inequality starts at $n=4$ and is preserved when
$n$ increases, as is verified after doubling the exponential term.
The description bound is therefore
\[
 B(n,s)\leq2^{n+(4+2n)s}.
\]
By definition $16(n+1)s\leq2^n$, and $4+2n\leq3(n+1)$ for $n\geq1$;
hence $(4+2n)s\leq3\cdot2^{n-4}$. Also
$n\leq2^{n-2}=4\cdot2^{n-4}$ for $n\geq4$, by a one-step induction.
Thus
\[
 n+(4+2n)s\leq7\cdot2^{n-4}\leq2^{n-1}.
\]
The previous theorem proves the stated count and constructs $h_n$.
The phrase ``almost all'' can consequently be given the exact finite
content that the fraction of tables with these small circuits is at most
$2^{-2^{n-1}}$. Only the cross-multiplied integer inequality is used in
\ipra.
\end{proof}

\begin{corollary}[R12: a primitive-recursive language outside polynomial circuits]
\label{thm:R12}
There is a primitive-recursive Boolean language $H$, decidable in
space $2^{O(n)}$, with no polynomial-size circuit family. For every
$c\geq1$ and $k\geq0$, the explicitly chosen length
\[
 n=2^{4(k+c+4)}
\]
has the property that no circuit with at most $c(n+1)^k$ gates computes
$H$ correctly on all inputs of that length. For each such circuit a least
mismatch input is primitive recursively obtainable.
\end{corollary}
\begin{proof}
Let $H(x)$ be the $x$th bit of $h_n$ on length $n\geq4$, with arbitrary
fixed values at the finitely many smaller lengths. Here $x$ is its
fixed-length binary value; the ambient word coding also records $n$.
The bounded construction of $h_n$ proves primitive recursiveness.

To verify the advertised length, set $u=k+c+4$ and $t=4u$. For $t\geq4$,
$2^t\geq t^2$, as follows by induction from $t=4$ and
$(t+1)^2\leq2t^2$. Moreover
$t^2>4+c+(k+1)(t+1)$, since the right side is at most
$4u^2+2u$ and $t^2=16u^2$. As $n=2^t$ and $n+1\leq2^{t+1}$,
\[
 16c(n+1)^{k+1}
 \leq2^{4+c+(k+1)(t+1)}<2^{2^t}=2^n.
\]
Therefore $c(n+1)^k\leq s(n)$. The defining property of $h_n$ excludes
all circuits up to that size. Scan the $2^n$ inputs to find a mismatch
with any specified candidate; the absence of a mismatch would contradict
the definition of the hard table.

For the space bound, store one candidate table of $2^n$ bits, one circuit
description of $O(s(n)\log(n+s(n)+2))=O(2^n)$ bits, the current circuit's
gate values, and binary loop counters. A brute-force comparison may
recompute every candidate circuit's value on each input, so no collection
of all truth tables is stored. The counters for the truth-table and
description enumerations themselves have $O(2^n)$ bits. The work space
is therefore $2^{O(n)}$; the time can be doubly exponential. This proves
$\mathrm{EXPSPACE}\not\subseteq\mathrm P/\mathrm{poly}$ by an explicit
recursive example. It asserts no reverse containment: unrestricted
nonuniform advice can encode nonrecursive languages.
\end{proof}

\begin{theorem}[R13: a constructive universal circuit upper bound]
\label{thm:R13}
Every truth table on $n\geq1$ inputs has a circuit with at most
$n2^n+n$ gates, constructed primitive recursively from the table.
For $n\geq64$ there is a compiler producing a circuit $C$ with
\[
 n\,\operatorname{size}(C)\leq18\cdot2^n.
\]
Together with Theorem~\ref{thm:R11}, this gives the order
$\Theta(2^n/n)$ for the maximum, over truth tables of length $2^n$, of
minimum circuit size. No real asymptotic expression is needed for its
integer formulation.
\end{theorem}
\begin{proof}
First precompute the $n$ negated inputs. For every input assignment with
table bit one, form the conjunction of its $n$ matching literals, then
OR these terms. At most $2^n$ terms are present, and the zero-term case
uses the constant zero wire. The bound $n2^n+n$ follows by charging
$n-1$ gates to each conjunction and at most one OR gate per term.
The correctness invariant says that each partial OR equals the union
of the already processed accepted assignments.

For the sharper bound put
$m=\lfloor\log_2 n\rfloor-1$ and $q=n-m$. Then
$n/4<2^m\leq n/2$. Split an input into a $q$-bit prefix $a$ and an
$m$-bit suffix $b$. Precompute, as separate circuits, all the
$2^{2^m}$ Boolean functions of $b$ by the preceding construction. Each
requires at most $m2^m+m\leq n^2$ gates, so this library costs at most
$n^2 2^{\lfloor n/2\rfloor}$ gates. For every prefix $a$, the input
truth table explicitly specifies one residual function $b\mapsto H(a,b)$,
whose output wire is now available.

Build all prefix indicators using a shared binary prefix tree. First
precompute the negations of the prefix input bits, at cost at most $n$.
From the constant-one empty indicator, form the indicators for both
extensions of every prefix by AND with the next literal. This costs
less than $2^{q+1}$ gates. For each full prefix, AND its indicator with
its residual-function wire, and OR the resulting $2^q$ terms. These last
operations cost at most $2^{q+1}$ gates. Exactly one indicator is true
on every input, so the output equals the specified residual function.
The total gate count is at most
\[
 n^2 2^{\lfloor n/2\rfloor}+n+4\cdot2^q.
\]
For $n\geq64$, $n^3\leq2^{\lfloor n/2\rfloor}$: check $64$ and $65$,
then use $(n+2)^3\leq2n^3$ for $n\geq8$ to advance by two.
Consequently the first term, multiplied by $n$, is at most $2^n$.
Also $n^2\leq2^n$. Finally $2^m>n/4$ gives
$n\cdot4\cdot2^q<16\cdot2^n$. Adding the three integer bounds gives
the claimed constant 18. All residual-table lookups and circuit wiring
are bounded operations, so this is a uniform primitive-recursive
compiler from a truth table to its circuit, although its input and output
can themselves be exponentially large in $n$.
\end{proof}

\subsection{Decision-tree and communication lower bounds}

A query tree is a coded finite rooted binary tree. Its internal vertices
specify which input bit to inspect, and its leaves carry output bits.
A comparison tree instead specifies comparisons between two input keys,
and its leaves give permutations of the keys. Repeated queries are
allowed; they can only increase the measured number of queries.

\begin{lemma}[R14: finite tree capacity]\label{thm:R14}
A binary tree of depth at most $h$ has at most $2^h$ leaves. An arbitrary
finite binary tree has at most $2^h-1$ leaves at depth strictly below $h$.
\end{lemma}
\begin{proof}
The first assertion follows by induction on $h$: a leaf contributes one,
and otherwise the two subtrees each contribute at most $2^{h-1}$.
Equivalently, map each leaf path to its extensions to length $h$; the
extension sets are nonempty and disjoint. For the second assertion,
encode a path of length $j<h$ by the integer $2^j$ plus its binary value.
There are only $1+2+\cdots+2^{h-1}=2^h-1$ such codes. Distinct leaves
give distinct path codes. These arguments can be implemented as bounded
counts over the tree code, so no induction on arbitrary set-valued trees
is required in the object language.
\end{proof}

\begin{theorem}[R15: comparison sorting requires many comparisons]
\label{thm:R15}
Any comparison tree that correctly sorts $n$ distinct keys has worst-case
depth $h$ satisfying
\[
 2^h\geq n!.
\]
If $n\geq4$, $a=\lfloor n/2\rfloor$, and
$\ell=\lfloor\log_2 a\rfloor$, then $h\geq a\ell$.
Under the uniform distribution on the $n!$ rank orders, the average
comparison count is also $\Omega(n\log n)$ in an explicit integer form:
for every $j$ with $2^j\leq n!/2$,
\[
 2\sum_{\pi}\operatorname{depth}(\pi)\geq j\,n!.
\]
\end{theorem}
\begin{proof}
Represent the $n!$ orders by permutations of ranks $0,\ldots,n-1$.
Each rank order has a unique permutation of positions that lists the
keys increasingly. Two distinct rank orders cannot reach the same leaf
of a correct tree, because that leaf specifies only one permutation.
The map from rank orders to leaves is therefore injective. Enumeration
of permutations, for example by factorial-radix codes, and a finite
occupancy count prove that there are at least $n!$ leaves. Apply
Lemma~\ref{thm:R14}.

The largest $a$ factors in $n!$ are each at least $a$, so
$n!\geq a^a\geq2^{a\ell}$. Strict monotonicity of powers of two gives
$h\geq a\ell$. This avoids Stirling's formula and real logarithms in
the proof. For the average bound, fewer than $2^j$ leaves can have depth
below $j$. Hence at least $n!/2$ of the $n!$ rank orders have depth at
least $j$, and summing their depths gives the displayed inequality.
Choosing $j=\lfloor\log_2(n!)\rfloor-1$ when this is nonnegative gives
the stated order bound. These conclusions are specifically about the
comparison model; integer-key algorithms using other operations are not
covered by the lower bound \cite{mit2010sorting}.
\end{proof}

\begin{theorem}[R16: exact deterministic query bounds]\label{thm:R16}
The deterministic query complexity of $n$-bit OR, AND, and parity is $n$.
For parity, every correctly answered input path must query all $n$ bits.
For OR the all-zero path must query all bits, and for AND the all-one
path must do so.
\end{theorem}
\begin{proof}
Suppose the all-zero OR path omits position $i$. Flip only that position.
Induction through the path shows every queried answer remains the same,
so the same leaf is reached, although the correct OR value changes.
The omitted position can be chosen as the least position absent from
the recorded query list. The AND argument starts from the all-one input.
For parity, flip any omitted position on any input; the transcript again
remains unchanged and parity changes. These are explicit error witnesses.
Reading all positions and folding OR, AND, or XOR gives matching upper
bounds. The empty-input case has zero queries with its stipulated
constant output.
\end{proof}

\begin{theorem}[R17: finite randomized query lower bounds]\label{thm:R17}
Suppose a randomized query algorithm uses a fixed finite uniform seed
space and at most $q$ queries on every input and seed, and errs on at
most one third of seeds for every input. For $n$-bit OR, $3q\geq n$.
For $n$-bit parity, $q\geq n$. The same parity conclusion applies to any
error bound strictly below one half.
\end{theorem}
\begin{proof}
Let the seed count be $R>0$. For OR, let $S_i$ be the set of seeds for
which the run on the all-zero input queries position $i$. Couple that
run to the run with only bit $i$ set, using the same seed. If the zero
run does not query $i$, induction on the transcript shows the two runs
are identical. They cannot both be correct because their OR values
differ. The two error sets each have size at most $R/3$, so
$R-|S_i|\leq2R/3$, or $3|S_i|\geq R$. On the other hand, each seed
occurs in at most $q$ of the sets $S_i$, whence
$\sum_i|S_i|\leq qR$. Summing the first inequalities gives
$nR\leq3qR$; cancel $R>0$.

For parity assume $q<n$. Fix a seed, obtaining a deterministic tree.
For each reachable leaf, choose the least input position not queried on
its path. Flipping that bit preserves the leaf and exchanges the two
parity values. It is an involution without fixed points on the inputs
reaching that leaf. Thus exactly half of those inputs are answered
correctly by the leaf's fixed output. Summing over leaves, each seed is
correct on exactly half of the $2^n$ inputs. Summing over the $R$ seeds
therefore gives exactly $R2^n/2$ correct input--seed pairs. But correctness
on at least two thirds of the seeds at every input gives at least
$2R2^n/3$ such pairs, a contradiction after multiplying by six. More
generally any strict majority guarantee contradicts the same exact
half-count. The argument uses finite double counting and an explicitly
computable pairing, not a minimax theorem.
\end{proof}

A deterministic two-party communication protocol is a finite binary tree.
Each internal vertex specifies the speaker and a Boolean function of
that speaker's input and the preceding transcript. A leaf gives the
output to both participants. Communication cost is the maximum transcript
length. Fixing this shared leaf-output convention matters for exact
additive constants.

\begin{lemma}[R18: rectangles and fooling sets]\label{thm:R18}
For each protocol leaf the input pairs reaching it form a Cartesian
product $X_0\times Y_0$. If a Boolean communication problem has a finite
set $F$ of one-inputs such that two distinct $(x,y),(x',y')\in F$ have
at least one zero cross-input $(x,y')$ or $(x',y)$, then every correct
protocol has at least $|F|$ one-leaves, and $2^h\geq|F|$ for its cost $h$.
\end{lemma}
\begin{proof}
At the root every input pair is possible. A message from Alice restricts
only the first coordinate of the current Cartesian product, and a
message from Bob restricts only the second. Induction on the transcript
proves the rectangle assertion. If two distinct elements of $F$ reached
one leaf, both cross-inputs would reach it as well. At least one has
value zero, contradicting the one label required by the two elements of
$F$. Thus $F$ injects into the one-leaves. Use the finite tree bound.
All sets can be represented by bounded Boolean matrices and row/column
bit vectors \cite{blais2025communication}.
\end{proof}

\begin{theorem}[R19: deterministic equality needs $n+1$ bits]
\label{thm:R19}
For $n\geq1$, equality of two $n$-bit strings has deterministic
communication complexity exactly $n+1$ under the shared leaf-output
convention.
\end{theorem}
\begin{proof}
The $2^n$ diagonal pairs $(x,x)$ form a fooling set: for $x\ne x'$, the
cross-pairs have unequal coordinates. Therefore at least $2^n$ leaves
must be labelled one. Since $n\geq1$, a zero-input exists and must reach
some zero-leaf. The tree has at least $2^n+1$ leaves, so $2^h>2^n$ and
$h\geq n+1$. For the upper bound Alice sends her $n$ input bits; Bob
compares with his input and sends the resulting answer bit. If only Bob
is required to know the final answer, this last bit can be omitted and
the corresponding exact bound becomes $n$. The $2^n$ diagonal pairs
still require distinct transcripts: a shared transcript for $(x,x)$ and
$(y,y)$ with $x\ne y$ would give Bob the same view on $(x,y)$ as on
$(y,y)$ and force a wrong answer. Thus the lower bound $n$ persists.
\end{proof}

\begin{theorem}[R20: a constant-communication public-randomness equality test]
\label{thm:R20}
With $kn$ shared random bits, equality has a protocol communicating
$k+1$ bits, having no false negatives and false-positive density exactly
$2^{-k}$ on unequal inputs. In particular three communicated bits suffice
for error at most one third, independently of $n$.
\end{theorem}
\begin{proof}
View the shared bits as $k$ independent masks $u_j\in\{0,1\}^n$.
Alice sends the $k$ parities $\langle x,u_j\rangle$ over the two-element
field. Bob compares each with $\langle y,u_j\rangle$ and sends one if
all match, zero otherwise. Equal strings always pass. For unequal strings
let $i$ be their least differing coordinate. Flipping bit $i$ of a mask
exchanges agreement and disagreement of the two inner products. Exactly
half the masks collide, by the pairing lemma. Independence, expressed
by the product count, gives exactly $2^{kn-k}$ colliding mask tuples.
Taking $k=2$ gives error $1/4$ and cost three. The bound uses public
randomness explicitly; it is not a claim about private-coin protocols.
\end{proof}

\subsection{Finite costs and amortized analysis}

\begin{theorem}[R21: best, weighted average, and worst cost]
\label{thm:R21}
Let $c_0,\ldots,c_{N-1}$ be natural costs with $N>0$, and let
$w_0,\ldots,w_{N-1}$ be natural weights of total $W>0$. If
$b=\min_i c_i$ and $z=\max_i c_i$, then
\[
 bW\leq\sum_{i<N}w_ic_i\leq zW.
\]
Both extrema and the weighted sum are primitive recursive on the coded
finite lists. These inequalities give best $\leq$ weighted average
$\leq$ worst after rational interpretation.
\end{theorem}
\begin{proof}
The bounded scan computing the minimum and maximum retains the least
index achieving the running extremum. For every $i$, multiplication of
$b\leq c_i\leq z$ by the nonnegative integer $w_i$ preserves order.
Add these inequalities and use $\sum_i w_i=W$. The proof is induction
on the list length with a numerical partial-sum invariant.
\end{proof}

There is no automatic ordering of an amortized bound between an
input-distribution average and a worst-case single-operation cost.
Amortized analysis controls totals along valid sequences of state changes;
it is not a third statistic of the same finite input distribution. The
following statement is the exact useful principle.

\begin{theorem}[R22: the potential method over natural numbers]
\label{thm:R22}
For a finite state sequence $s_0,\ldots,s_m$, let $c_i$ be the cost of
the transition $s_i\to s_{i+1}$, and let $\Phi(s)$ be a nonnegative
integer potential. If
\[
 c_i+\Phi(s_{i+1})\leq a_i+\Phi(s_i)\qquad(i<m),
\]
then
\[
 \sum_{i<m}c_i+\Phi(s_m)
 \leq\sum_{i<m}a_i+\Phi(s_0).
\]
In particular $\sum_{i<m}c_i\leq\sum_{i<m}a_i+\Phi(s_0)$.
\end{theorem}
\begin{proof}
For $m=0$ the inequality is equality. Add the next transition inequality
to the induction hypothesis and cancel the common intermediate potential
$\Phi(s_m)$. This gives the inequality for $m+1$. The final consequence
uses $\Phi(s_m)\geq0$. No subtraction in the natural numbers, and no
assumption that each amortized charge is a realized single-operation
cost, is needed.
\end{proof}

\begin{corollary}[R23: binary counters have constant amortized bit-flip cost]
\label{thm:R23}
Starting at zero, $m$ increments of an unbounded binary counter use at
most $2m$ bit flips. More precisely, if $\nu(m)$ is the number of one bits
in the final counter, the total number of flips is $2m-\nu(m)$.
\end{corollary}
\begin{proof}
If the current value has $t$ trailing ones, the next increment changes
those $t$ ones to zero and the following zero to one. Its cost is $t+1$,
and the potential $\Phi=\nu$ changes from $\nu$ to $\nu-t+1$.
Thus $c_i+\Phi(s_{i+1})=2+\Phi(s_i)$ exactly. Sum this identity and use
$\Phi(0)=0$ to obtain total cost plus $\nu(m)$ equal to $2m$.
The trailing-one count, population count, and binary increment are all
bounded digit scans. This is a bit-flip statement: a particular tape
implementation may require additional head-motion cost, which must be
accounted for separately.
\end{proof}


\subsection{Lean targets and the remaining proof-lowering work}

The file \texttt{lean/CountingOutline.lean} defines the semantic interfaces
and target propositions for this section. It is an \emph{uncompiled outline},
with no axioms or admitted theorem bodies, rather than a completed proof of
these results. Its two elementary recursion equations are the only theorem
bodies. The substantive targets are definitions of propositions whose proofs
remain to be implemented.

The first layer defines \texttt{sumBelow}, \texttt{countBelow}, least-witness
specifications, product and union count targets, and finite random-machine
families. The proof order is to establish the two recursion equations,
induct on loop bounds for addition and monotonicity, prove radix product
factorization, and then implement first-witness search. The weighted identity
in R03 uses the same product lemma with values one and two in place of bits.
R07 and R08 then require only the finite union bound, XOR permutations,
and the machine-to-circuit compiler already developed in the NP section.
The zero-error interface explicitly records truncated cost and answer;
its operational instantiation must prove consistency of prefixes,
absorbing halting states, and the fact that an unfinished $T$-step trace
has cost exactly $T$. These conditions cannot be replaced by unrelated
semantic cost and answer functions.

The second layer defines gates, topologically valid circuits, a circuit
interpreter, truth tables, the description bound, and hard-table and
universal-upper-bound targets. The implementations follow the proof:
bounded description enumeration, a marking-count invariant, and least
missing-table search; then literal conjunctions, prefix indicators, and
residual-function selection. The final layer supplies finite query trees,
comparison trees, and communication protocols, together with their
interpreters and depth functions. The principal lemmas are preservation of
transcripts under an unqueried bit flip and preservation of Cartesian
products along protocol paths. Counts of leaves, permutations, and mask
pairs finish their lower bounds.

All these Lean function and inductive-type parameters are a semantic
presentation of the mathematics. For an actual \ipra{} derivation, instantiate
them with the fixed primitive-recursive programs, encode every finite object
as a natural number, and replace every bounded property by its Boolean
compiler. Prove the loop invariants with the repository's quantifier-free
induction rule and check the resulting object-language derivation. A kernel
proof of an ambient Lean target alone would not establish that last fact.
The thirteen independent checks in \texttt{code/random\_lower\_checks.py}
exercise finite instances and exact identities; they supply implementation
evidence and do not replace universal proofs.
